\documentclass{article}
\usepackage[T1]{fontenc}
\usepackage{lmodern}
\usepackage{graphicx}
\usepackage{amsmath,amssymb}
\usepackage[a4paper,margin=2cm]{geometry}
\usepackage[absolute,overlay]{textpos}
\setlength{\TPHorizModule}{1mm}
\setlength{\TPVertModule}{1mm}
\usepackage[indonesian]{babel}
\usepackage{hyperref}
\setlength{\emergencystretch}{2em}
% unit: Halaman 7

\begin{document}

% === Halaman 7 (llm) ===

\begin{itemize}
    \item Sebuah grup $\langle G, *\rangle$ dikatakan \textcolor{red}{grup komutatif} atau \textcolor{red}{grup abelian} (atau disingkat \textcolor{red}{abelian} saja) jika berlaku aksioma komutatif $a * b = b * a$ untuk semua $a, b \in G$.
    \item $\langle R, +\rangle$ dan $\langle R, \cdot\rangle$ adalah abelian
    \item $\langle Z, +\rangle$ dan $\langle Z, \cdot\rangle$ adalah abelian
    \item tetapi, $\langle M, \times\rangle$, dengan $M$ adalah himpunan matriks $2 \times 2$ dengan determinan $\neq 0$ bukan abelian (tanya kenapa?)
\end{itemize}

\vspace{10mm}

Ket: Abelian diambil dari kata ``abel'', untuk menghormati Niels Abel, seorang Matematikawan Norwegia (1802 -- 1829)

\end{document}
